\documentclass{article}
\usepackage[T1]{fontenc}
\usepackage[utf8]{inputenc}
\usepackage{amsmath, amssymb}
\usepackage{hyperref}
\usepackage{bookmark}
\usepackage{xurl}
\usepackage{longtable}
\usepackage{siunitx} % For better handling of numbers and units
\usepackage{calc}
\usepackage{array}
\usepackage{booktabs} % For better table formatting
\usepackage{parskip} % For better handling of paragraph spacing
\usepackage[a4paper, margin=2.5cm]{geometry} % Adjust margins
% Définition de \tightlist
\providecommand{\tightlist}{%
  \setlength{\itemsep}{0pt}\setlength{\parskip}{0pt}}

% Assouplit la justification
\sloppy 
\setlength{\parindent}{0pt} % Pas d'indentation

\title{Logistic Regression: Introduction and Interpretation}

\author{David Thébault}
\date{October 18, 2025}


\begin{document}

\maketitle % Generate the title

\section{\texorpdfstring{\textbf{Logistic
Regression}}{Logistic Regression}}\label{logistic-regression}

\textbf{Objective}: Understand, implement, and apply logistic regression
for diabetes prediction.


\subsection{Introduction to Logistic Regression}

Logistic regression applies when the dependent variable is binary.
Formally, let \(Y\) be a random variable taking values in \(\{0,1\}\),
where \(Y=1\) indicates that the event occurs and \(Y=0\) otherwise.
Let \(X_1,\ldots,X_k\) denote a set of \(k\) explanatory variables,
which may be either qualitative or quantitative.

We consider the joint population vector \((Y, X_1,\ldots,X_k)\).
A random sample of size \(n\) is drawn from this population,
yielding observations \((Y_i, X_{i1},\ldots,X_{ik})\) for
\(i = 1,\ldots,n\).

Unlike simple linear regression, logistic regression 
\textbf{estimates the probability} of an event occurring, rather than predicting a
specific numerical value.

\subsection{Theory and Fundamentals}

\subsubsection{Probabilistic Model}

This section describes the underlying probabilistic framework of the
model, focusing on the random variable and its distribution.

The variable \(Y_i\) follows a \textbf{Bernoulli distribution} with
parameter \(p_i\) representing the probability that \(Y_i=1\).

\[Y_i \sim B(p_i)\]

The probabilities are defined as:

\[P(Y_i=1) = p_i \quad, \quad P(Y_i = 0) = 1 - p_i\]

This can also be written as::

\[P(Y_i = k) = {p_i}^k(1 - p_i)^{1-k} \quad \text{for k} \in \{0, 1\}\]

\subsubsection{LOGIT Model}\label{logit-model}

This section focuses on the specific form of the logistic regression
model, which is a type of probabilistic model.

\medskip

To ensure that the expected value of \(Y, E(Y)\), only takes values
between 0 and 1, we use the \textbf{logistic function}:

\[f(x) = \dfrac{\text{exp(x)}}{1 + \text{exp(x)}} = p\]

or equivalently:

\[f(x) = \dfrac{1}{1 + \text{exp(-x)}} = p\]

\medskip

This guarantees that \(0 < f(x) < 1\), so \(E[Y]\) can represent a valid
probability.

The \textbf{logit function} transforms a probability \(p\) into an
\textbf{unrestricted real value}:

\textbf{Notations}:

\begin{itemize}
\tightlist
\item
  \(X = (1,X_1, \ldots, X_k)\) vector of predictors, including the
  intercept term.
\item
  \(\beta = (\beta_0,\beta_1, \ldots, \beta_k)\) vector of coefficients.
\end{itemize}

\medskip

\textbf{Logit Transformation}

The logit of a probability \(p\) is defined as:

\[\text{logit}(p) = \text{log}(\dfrac{p}{1 - p})\]

In the context of logistic regression, the logit of the probability is
modeled as a linear combination of the predictors:

\[\text{logit}(p) = \beta .X\]

Explicitly, this can be written as:

\[\text{logit}(p) = \beta_0 + \beta_1 X_1 + \beta_2 X_2 + \ldots + \beta_k X_k\]

Or equivalently:

\[\text{log}\left( \dfrac{p}{1-p} \right) = \beta_0 + \beta_1 X_1 + \beta_2 X_2 + \ldots + \beta_k X_k\]

\medskip

\textbf{Logistic Function}(inverse of logit).

To convert the linear predictor back into a probability, we use
the logistic function which is the inverse of the logit function:

\[p = \frac{1}{1 + \exp(-\beta .X)}\]

Demonstration:

\[p(x) = \dfrac{1}{1 + \exp(-\beta x)}\]

\[\underset{inverse}   \iff \dfrac{1}{p} = 1 + \exp(-\beta x)\]

\[\iff \dfrac{1}{p} - 1 = \exp(-\beta x)\]

\[\iff \dfrac{1}{p} - \dfrac{p}{p} = \exp(-\beta x)\]

\[\iff \dfrac{1-p}{p} = \exp(-\beta x)\]

\[\iff \log(\dfrac{1-p}{p}) = -\beta x\]

\[\iff \log(\dfrac{p}{1-p}) = \beta x\]

For simplicity, we often write \(p\) instead of \(p(x)\) when the
context is clear.

\subsubsection{Model Assumptions}

For the logistic regression model to be valid and
generalizable, the following key assumptions must be satisfied.
Violations of these assumptions can lead to biased or inefficient
estimates, and may compromise the interpretability of the model.

\medskip

\textbf{I. Linearity of Log-Odds}

\textbf{The relationship between each continuous predictor (\(X_j\)) 
and the log-odds of the outcome (\(Y=1\))
must be linear}. 

Mathematically, this means:
\[\log\left(\frac{P(Y=1|X)}{1 - P(Y=1|X)}\right) = \beta_0 + \beta_1 X_1 + \beta_2 X_2 + \cdots + \beta_k X_k\]

Implications:

\begin{itemize}
\tightlist
\item
  If the true relationship is non-linear (e.g., quadratic or U-shaped),
  the linear assumption will not hold.
\item
  How to check: Plot the log-odds against each continuous
  predictor. If the relationship appears curved, consider adding
  polynomial terms or splines.
\end{itemize}

Example: If \texttt{age} is a predictor, the log-odds of
diabetes should increase or decrease linearly with \texttt{age}. If the
relationship is non-linear, you might need to add a quadratic term
(e.g., \texttt{age²}).

\medskip

\textbf{II. No Multicollinearity}

\textbf{Predictor variables should not be highly correlated with each other}. 

High multicollinearity can inflate the variance of coefficient estimates, 
making them unstable and difficult to interpret.

Implications:

\begin{itemize}
\tightlist
\item
  Multicollinearity does not bias the model, but it reduces the
  precision of the estimated coefficients.
\item
  How to check: For numeric variables, use Variance Inflation Factor (VIF) 
  (see multicolinearity assessment). A   VIF \textgreater{} 5 or 10 
  indicates problematic multicollinearity.
\end{itemize}

Example: If your model includes both \texttt{blood pressure}
and \texttt{blood pressure squared}, these variables are likely to be
highly correlated. Consider removing one or using regularization (e.g.,
Lasso or Ridge).

\medskip

\textbf{III. Additivity}

\textbf{The effect of each predictor on the log-odds of the
outcome is additive}. 

This means the contribution of each predictor to the log-odds 
is independent of the values of the other predictors.

Implications:

\begin{itemize}
\tightlist
\item
  If there are interaction effects (e.g., the effect of
  \texttt{age} on diabetes depends on \texttt{BMI}), they must be
  explicitly modeled.
\item
  How to check: Test for interactions by adding product terms
  (e.g., \texttt{age x BMI}) to the model and evaluating their
  significance.
\end{itemize}

Example: If the effect of \texttt{glucose} on diabetes risk
depends on \texttt{age}, you should include an interaction term like
\texttt{glucose age} in your model.

\medskip

\textbf{IV. Independence of Observations}

\textbf{The model assumes that observations are independent of each other}.

This is critical for the validity of standard errors and hypothesis tests.

Implications:

\begin{itemize}
\tightlist
\item
  If observations are correlated (e.g., repeated measurements from the
  same individual or clustered data), standard errors will be
  underestimated, leading to inflated Type I error rates.
\item
  How to check: Look for patterns in residuals or use tests for
  autocorrelation (e.g., Durbin-Watson test for time-series data).
\end{itemize}

Example: If your dataset includes multiple measurements from
the same patient over time, you should use a mixed-effects
logistic regression or GEE (Generalized Estimating Equations)
to account for within-patient correlation.

\subsection{Practical Implementation of Logistic}

\paragraph{Recommendations}

\begin{enumerate}
\tightlist
\item
  \textbf{Always check assumptions} before interpreting your model.
\item
  \textbf{Use diagnostic plots} (e.g., residual plots, VIF) to identify
  violations.
\item
  \textbf{Consider regularization} (Lasso/Ridge) if multicollinearity is
  an issue.
\item
  \textbf{Model interactions} if additivity is violated.
\item
  \textbf{Use robust standard errors} or mixed-effects models if
  observations are not independent.
\end{enumerate}


\paragraph{\texorpdfstring{\textbf{Fit the logistic regression
model}}{Fit the logistic regression model}}\label{fit-the-logistic-regression-model}

\subparagraph{Logistic Regression with
scipy}\label{logistic-regression-with-scipy}

\paragraph{\texorpdfstring{\textbf{Linearity of
Log-Odds}}{Linearity of Log-Odds}}\label{linearity-of-log-odds}

Predicted Probabilities and Log-Odds:

\begin{itemize}
\tightlist
\item
  The linearity of log-odds assumption in logistic regression refers to
  the relationship between continuous predictors and the log-odds of the
  outcome.
\item
  To assess this, you need the predicted probabilities from a fitted
  logistic regression model. These probabilities are then transformed
  into log-odds.
\end{itemize}

Visual Inspection:

\begin{itemize}
\tightlist
\item
  You plot the continuous predictors against the predicted log-odds to
  visually inspect if the relationship is linear.
\item
  If the relationship is not linear, you may need to transform the
  predictor (e.g., using polynomials, splines, or categorization).
\end{itemize}

\subparagraph{Calculate Predicted Probabilities and
Log-Odds}\label{calculate-predicted-probabilities-and-log-odds}

\subparagraph{Plot Log-Odds vs. Continuous
Predictors}\label{plot-log-odds-vs-continuous-predictors}

\subsection{Coefficients
interpretation}\label{coefficients-interpretation}

\textbf{The Odds}

The odds are defined by:

\[\text{Odds} = \dfrac{p}{1-p}\]

\(\text{Where} \quad p = P(target=1|X)\)

\begin{quote}
\emph{If a student has a 3 in 4 chance of passing and a 1 in 4 chance of
failing, their odds are \textquotesingle3 to 1\textquotesingle:}
\(\text{Odds} = \dfrac{3/4}{1/4}=3\)
\end{quote}

\begin{itemize}
\tightlist
\item
  \textbf{Notation:}
  \[\text{Odds}(Y=1|X=0)=\dfrac{P(Y=1|X=0)}{1-P(Y=1|X=0)}\]
\end{itemize}

\subsubsection{\texorpdfstring{\textbf{The Odds
Ratio}}{The Odds Ratio}}\label{the-odds-ratio}

The odds ratio comparing the \textbf{probability of \(target=1\)}
between individuals with value \(X\) and those without it.

\[\text{Odds Ratio} = \dfrac{\text{Odds}(Y=1|X=1)}{\text{Odds}(Y=1|X=0)}\]

\[\text{Odds Ratio} = \dfrac{P(Y_i=1 | X=1)}{1 - P(Y_i=1 | X=1)} / \dfrac{P(Y_i=1 | X=0)}{1 - P(Y_i=1 | X=0)}\]

We know that logit is given by:

\[\text{logit}(p) = \text{log}(\dfrac{p}{1-p}) = \beta_0 + \beta_1 x_{i1} + \beta_2 x_{i2} + \ldots + \beta_k x_{ik}\]

\subsubsection{\texorpdfstring{\textbf{Interpreting the
Intercept}}{Interpreting the Intercept}}\label{interpreting-the-intercept}

The intercept \(\beta_0\)\hspace{0pt} represents the \textbf{log-odds of
the outcome \(Y=1\) when all predictors are equal to zero}.\\
\(\beta_0\)\hspace{0pt} defines the \textbf{baseline probability} of the
outcome when all predictors are zero.

\textbf{Caveat}:\\
This interpretation of \(\beta_0\) is often not meaningful if some
predictors cannot logically be zero (e.g., age=0, blood pressure). In
such cases, \(\beta_0\)\hspace{0pt} is primarily a mathematical
component of the model and is rarely interpreted in isolation.

\begin{itemize}
\tightlist
\item
  \textbf{Odds for the baseline group:}
\end{itemize}

\[\text{Odds}(Y=1 | X_1=0) = \exp(\beta_0)\]
\begin{itemize}
\tightlist
\item
  \textbf{Probability for the baseline group:}
  \[P(Y=1 | X_1=0)=\dfrac{\exp(\beta_0)}{1 + \exp(\beta_0)}\]
\end{itemize}

\begin{quote}
If \(X_1\)\hspace{0pt} is "smoking status" (\(0\) = non-smoker, \(1\) =
smoker), then

\begin{itemize}
\tightlist
\item
  \(\beta_0\)\hspace{0pt} gives the \textbf{log-odds} of the outcome for
  non-smokers\\
\item
  \(\exp(\beta_0)\) gives the \textbf{odds} of the outcome for
  non-smokers.
\end{itemize}

If \(\beta_0 = -1\), then: \[\exp(\beta_0) = \exp(-1) \approx 0.37\]

\[P(Y=1 | X_1=0)= \dfrac{0.37}{1 + 0.37} \approx 0.27\]

27% of non-smokers are predicted to have the outcome (e.g., lung
cancer), assuming no other predictors.\\
It is the observed proportion of lung cancer for non-smokers.
\end{quote}

\subsubsection{\texorpdfstring{\textbf{Interpreting the
Slope}}{Interpreting the Slope}}\label{interpreting-the-slope}

In a model with multiple predictors, each \(\beta_i\)\hspace{0pt} (and
its corresponding odds ratio \(\exp(\beta_i)\) represents the effect of
that predictor on the log-odds of \(Y=1\), holding all other predictors
constant.\\
This is the key assumption of multivariable regression: ceteris paribus
(all else being equal).

The coefficient \(\beta_1\)\hspace{0pt} represents the \textbf{change in
the log-odds of} \(Y=1\) for a \textbf{one-unit change} in
\(X_1\)\hspace{0pt}. The odds ratio \(\exp(\beta_1)\)\hspace{0pt}
quantifies how the odds of \(Y=1\) change with \(X_1\)\hspace{0pt}.

\textbf{General Formula for Odds Ratio}

For any type of predictor \(X_1\), the odds ratio for a one-unit
increase is:\\
\[\text{Odds Ratio} = \frac{\text{Odds}(Y=1 | X_1 = x+1)}{\text{Odds}(Y=1 | X_1 = x)} = \exp(\beta_1)\]

\textbf{Note:}\\
\(\exp(\beta_1)\) compares the odds of \(Y=1\) between \(X_1=1\) and
\(X_1=0\), controlling for all other variables in the model (all others
features constant).

\begin{itemize}
\tightlist
\item
  Case: \(X_1\) is Binary
\end{itemize}

For a binary predictor \(X_1\)\hspace{0pt} (e.g., \(0\) = non-smoker,
\(1\) = smoker), the odds ratio \(\exp(\beta_1)\)\hspace{0pt} compares
the odds of \(Y=1\) between the two groups.

\begin{itemize}
\tightlist
\item
  \textbf{Logistic regression equation:}
\end{itemize}

\[\log\left(\dfrac{P(Y=1 | X_1)}{1 - P(Y=1 | X_1)}\right) = \beta_0 + \beta_1 1_{\{X_1 = 1\}}\]

\begin{itemize}
\tightlist
\item
  \textbf{Odds ratio:}
  \[\text{Odds Ratio} = \dfrac{P(Y=1 | X_1=1)}{1 - P(Y=1 | X_1=1)} / \dfrac{P(Y=1 | X_1=0)}{1 - P(Y=1 | X_1=0)} = \exp(\beta_1)\]
\end{itemize}

\textbf{Interpretation:}

\begin{itemize}
\tightlist
\item
  If \(\exp(\beta_1) = 1\): No effect of the feature \(X_1\)\hspace{0pt}
  on the odds of \(Y=1\).
\item
  If \(\exp(\beta_1)>1\): The odds of \(Y=1\) are higher when \(X_1=1\).
  The feature \(X_1\)\hspace{0pt} is \textbf{positively associated} with
  the outcome.
\item
  If \(\exp(\beta_1) < 1\): The odds of \(Y=1\) are lower when
  \(X_1=1\). The feature \(X_1\)\hspace{0pt} is \textbf{negatively associated} with the outcome.
\end{itemize}

\begin{quote}
\textbf{Example:}\\
If \(\beta_1 = 0.7 \rightarrow \exp(\beta_1) \approx 2.01\). The odds of
lung cancer for smokers \((X_1=1)\) are twice as high as for non-smokers
\((X_1=0)\).
\end{quote}

\begin{itemize}
\tightlist
\item
  Case: \(X_1\) is Categorical
\end{itemize}

For a categorical predictor \(X_1\) with more than two levels (e.g.,
color = red, green, blue), you use \textbf{dummy variables}.

\begin{itemize}
\tightlist
\item
  \textbf{The logistic regression model becomes:}
\end{itemize}

\[\text{log}\left( \dfrac{P(Y=1)}{1 - P(Y=1)} \right) = \beta_0 + \beta_{green}1_{\{X_1 = \text{green}\}} + \beta_{blue}1_{\{X_1 = \text{blue}\}}\]

\begin{itemize}
\tightlist
\item
  \textbf{Reference Category ("red"):} When
  \(1_{\{X_1 = \text{green}\}}=0\) and \(1_{\{X_1 = \text{blue}\}}=0\),
  the log-odds are:
\end{itemize}

\[\text{log}\left( \dfrac{P(Y=1)}{1 - P(Y=1)} \right) = \beta_0\]

\textbf{Interpretation:} This means \(\beta_0\)\hspace{0pt} represents
the log-odds of \(Y=1\) for the reference category ("red").

\begin{itemize}
\tightlist
\item
  \textbf{Category ("green"):} When \(1_{\{X_1 = \text{green}\}}=1\) and
  \(1_{\{X_1 = \text{blue}\}}=0\), the log-odds are:
\end{itemize}

\[\text{log}\left( \dfrac{P(Y=1)}{1 - P(Y=1)} \right) = \beta_0 + \beta_{green}\]

\begin{itemize}
\tightlist
\item
  \textbf{Category ("blue"):} When \(1_{\{X_1 = \text{green}\}}=0\) and
  \(1_{\{X_1 = \text{blue}\}}=1\), the log-odds are:
\end{itemize}

\[\text{log}\left( \dfrac{P(Y=1)}{1 - P(Y=1)} \right) = \beta_0 + \beta_{blue}\]

The \textbf{odds ratio for "blue" relative to the reference "red"} is:
\[\exp(\beta_{blue}) = \dfrac{\text{Odds}(Y=1 | \text{blue})}{\text{Odds}(Y=1 | \text{red})}\]

The same way, \(\exp(\beta_{\text{green}})\)\hspace{0pt} compares the
odds for "green" vs. the "red" reference.

\begin{quote}
\textbf{Interpretation:}

If \(\exp(\beta_{\text{green}}) = 1.5\), the odds of \(Y=1\) are \(1.5\)
times higher for "green" compared to "red".
\end{quote}

\begin{itemize}
\item ~
  \subparagraph{\texorpdfstring{Case: \(X_1\) is
  Quantitative}{Case: X_1 is Quantitative}}\label{case-x_1-is-quantitative}
\end{itemize}

For a continuous predictor \(X_1\) (e.g., age, blood pressure), the odds
ratio \(\exp(\beta_1)\)\hspace{0pt} represents the multiplicative change
in the odds of \(Y=1\) for a one-unit increase in \(X_1\)\hspace{0pt}.

\begin{itemize}
\tightlist
\item
  \textbf{Logistic regression equation:}
\end{itemize}

\[\log\left(\dfrac{P(Y=1 | X_1)}{1 - P(Y=1 | X_1)}\right) = \beta_0 + \beta_1 X_1\]

\begin{itemize}
\tightlist
\item
  \textbf{Odds ratio for a one-unit increase:}
\end{itemize}

\[\text{Odds Ratio} = \frac{\text{Odds}(Y=1 | X_1 = x+1)}{\text{Odds}(Y=1 | X_1 = x)} = \exp(\beta_1)\]

\textbf{In short}: \(\beta_1\)\hspace{0pt} captures the \textbf{constant
log-odds} change per unit increase in \(X_1\), so
\(\exp(\beta_1)\)\hspace{0pt} is the \textbf{odds ratio} for that
one-unit change.

This holds regardless of the starting value of \(X_1\)\hspace{0pt}
because the model assumes a constant multiplicative effect on the odds
(a key assumption of logistic regression).

\begin{quote}
\textbf{Interpretation:}

\begin{itemize}
\tightlist
\item
  If \(\beta_1=0.095 \rightarrow \exp(\beta_1)=1.1\), the odds of
  \(Y=1\) increase by 10% for each one-unit increase in
  \(X_1\)\hspace{0pt}.
\item
  If \(X_1\)\hspace{0pt} is "years of smoking" and
  \(\beta_1 = 0.7 \rightarrow  \exp(\beta_1) \approx 2.01\). For each
  additional year of smoking, the odds of lung cancer double. .
\end{itemize}
\end{quote}

\textbf{Summary}

\begin{longtable}[]{@{}
  >{\raggedright\arraybackslash}p{(\columnwidth - 2\tabcolsep) * \real{0.2396}}
  >{\raggedright\arraybackslash}p{(\columnwidth - 2\tabcolsep) * \real{0.7604}}@{}}
\toprule\noalign{}
\begin{minipage}[b]{\linewidth}\raggedright
Type of \(X_1\)
\end{minipage} & \begin{minipage}[b]{\linewidth}\raggedright
Interpretation of \(\exp(\beta_1)\)
\end{minipage} \\
\midrule\noalign{}
\endhead
\bottomrule\noalign{}
\endlastfoot
Binary & Compares odds of \(Y=1\) between \(X_1=1\) and \(X_1=0\) \\
Categorical & Compares odds of \(Y=1\) for a given category relative to
the reference. \\
Quantitative & Multiplicative change in odds of \(Y=1\) for a one-unit
increase in \(X_1\) \\
\end{longtable}

\begin{enumerate}
\tightlist
\item
  Interpretation of the Plots
\end{enumerate}

\begin{itemize}
\item
  Consistency of Log-Odds: The boxplot will show the distribution of
  predicted log-odds for each category. If the medians of the log-odds
  for each category are distinct, it suggests that the categorical
  predictor is having an effect on the log-odds of the outcome.
\item
  Effect of Categories: The difference in log-odds between categories
  should reflect the coefficients in the logistic regression output.
\end{itemize}

\textbf{Interpreting the coefficients:}

For each feature, the exponentiated coefficient (exp(coef)) represents
the change in odds for a one-unit increase in that feature, holding all
other features constant.

For example, has the coefficient for
\textquotesingle pregnancies\textquotesingle{} is 0.1604, the odds ratio
is exp(0.1604) = 1.174037.\\
This means that for each one-unit increase in pregnancies, the odds of
the outcome occurring (e.g., having diabetes) increase by approximately
17.4%, assuming all other features remain constant.

\textbf{Diabetes Prediction: data preparation and model inference}

\begin{itemize}
\tightlist
\item
  Preparing new input data for prediction
\end{itemize}

\begin{itemize}
\tightlist
\item
  Using the model to make a prediction
\end{itemize}

\begin{itemize}
\tightlist
\item
  Transform continuous to categorical
\end{itemize}

\subsection{APPENDIX}\label{appendix}

\subsubsection{Interpreting the
coefficients}\label{interpreting-the-coefficients}

\textbf{Demonstration:} The coefficient \(\beta_1\)\hspace{0pt}
represents the \textbf{change in the log-odds of} \(Y=1\) for a
\textbf{one-unit change} in \(X_1\)\hspace{0pt} quantitative feature.

Notations:\\
\[\text{Odds}(Y=1|X=x+1)=P(Y=1|X=x+1) / (1-P(Y=1|X=x+1))\]

\[\text{Odds}(Y=1|X=x)=P(Y=1|X=x) / (1-P(Y=1|X=x))\]

We know that:

\[\text{log(Odds)}(Y=1|X=x+1)=\beta_0 + \beta_1 \times (x+1)\]

\[\text{log(Odds)}(Y=1|X=x)=\beta_0 + \beta_1 \times x\]

By difference:

\[\text{log(Odds)}(Y=1|X=x+1) - \text{log(Odds)}(Y=1|X=x) =\beta_0 + \beta_1 \times (x+1) - (\beta_0 + \beta_1 \times x) = \beta_1\]

\[\text{log}\left(\dfrac{\text{Odds}(Y=1|X=x+1)}{\text{Odds}(Y=1|X=x)}\right) =\beta_1\]

\textbf{CQFD}

Note:

\[\dfrac{\text{Odds}(Y=1|X=x+1)}{\text{Odds}(Y=1|X=x)} = \exp(\beta_1)\]

\subsubsection{Model formulation}\label{model-formulation}

The prediction \(y_i=1\) of the logistic regression is defined:

\[\hat{y_i} = P(y_i=1 | x_i; \theta) = \frac{1}{1 + \exp(-\theta^Tx_i)} = h_{\theta}(x_i)\]

\begin{itemize}
\tightlist
\item
  If \(y_i=1\), then \(P(y_i|x_i; \theta)=P(y_i=1|x_i; \theta)\)
\item
  If \(y_i=0\), then
  \(P(y_i|x_i; \theta)=P(y_i=0|x_i; \theta) = 1 - P(y_i=1|x_i; \theta)\)
\end{itemize}

We can write these two equations into a single one:

\[P(y_i|x_i; \theta)=P(y_i=1|x_i; \theta)^{y_i}\times (1 - P(y_i=1|x_i; \theta))^{1-y_i}\]

With the notations:

\[P(y_i|x_i; \theta)=h_{\theta}(x_i)^{y_i}\times (1 - h_{\theta}(x_i))^{1-y_i}\]

\subsubsection{Likelihood function}\label{likelihood-function}

The \textbf{likelihood} of the observations \(y_i\) given the inputs
\(x_i\) and parameters \(\theta\) is defined as:

\[L(\theta) = \prod_{i=1}^n P(y_i|x_i;\theta) = \prod_{i=1}^n (h_{\theta}(x_i))^{y_i} (1 -h_{\theta}(x_i))^{1-y_i}\]

where the prediction of the logistic regression is defined:

\[h_{\theta}(x_i) = P(y_i=1 | x_i; \theta) = \frac{1}{1 + \exp(-\theta^Tx_i)}\]

The \textbf{log-likelihood} is defined as:

\[l(\theta) = \log(L(\theta))=\sum_{i=1}^n[y_i \log (h_{\theta}(x_i)) + (1 - y_i) \log(1 - h_{\theta}(x_i))]\]

\subsubsection{Objective of Logistic
Regession}\label{objective-of-logistic-regession}

The goal of learning a \textbf{logistic regression model} is to
\textbf{minimize the cost function} by adjusting the parameters
\(\theta\).The cost function measures the average prediction error
across all \(n\) training samples.

\subsubsection{Cost function (general
definition)}\label{cost-function-general-definition}

The cost function \(J(\theta)\) is defined as the average penalty for
prediction errors across the training set. Mathematically, it is
expressed as:

\[J(\theta) = -\frac{1}{n} \sum cost(h_\theta(x_i), y_i)\]

where:

\begin{itemize}
\tightlist
\item
  \(h_\theta(x_i)\) the model\textquotesingle s prediction for sample
  \(x_i\)
\item
  \(y_i\) the observed (true) value.
\end{itemize}

Alternatively, it can be written as:

\[J(\theta) = \frac{1}{n} \sum cost(\hat{y_i}, y_i)\]

Where:

\begin{itemize}
\tightlist
\item
  \(\hat{y_i} = h_{\theta}(x_i)\)
\end{itemize}

\subsubsection{\texorpdfstring{Cost function \textbf{log
loss}}{Cost function log loss}}\label{cost-function-log-loss}

For logistic regression, the cost function is called \textbf{log loss}
(or logistic loss).\\
\textbf{Log loss} is derived from the log-likelihood \(l(\theta)\) and
is defined as:

\[J(\theta) = -\dfrac{1}{n}l(\theta) = -\frac{1}{n} \sum(y_i \log(h_\theta(x_i)) + (1 - y_i) \log(1 - h_\theta(x_i)))\]

\paragraph{How Log-Loss penalizes prediction
errors}\label{how-log-loss-penalizes-prediction-errors}

The log-loss function penalizes prediction errors based on the estimated
probability from the model. It assigns higher penalties when predictions
are far from the true labels---specifically:

\begin{itemize}
\item
  If the label \(y_i = 1\) (positive class), the penalty is
  \(-log(h_{\theta}(x_i))\). The closer \(h_{\theta}(x_i)\) is to 0 (far
  from the true label), the higher the penalty
  (\(-\log(0^+) \approx + \infty\)).
\item
  If the label \(y_i = 0\) (negative class), the penalty is
  \(-log(1-h_{\theta}(x_i))\), the closer \(h_{\theta}(x_i)\) is to 1
  (far from the true label), the higher the penalty
  (\(-\log(1 -1^+) \approx + \infty\)).
\end{itemize}

The two cases are combined into a single formula for observation \(i\):
\[y_i log(h_\theta(x_i)) + (1 - y_i) log(1 - h_\theta(x_i))\]

\textbf{Key insights:}

\begin{itemize}
\tightlist
\item
  \textbf{Log loss} evaluates how well the model fits the training data.
\item
  The log loss is the negative average the log likelihood.
\item
  \textbf{Higher likelihood} leads to \textbf{lower log loss} (since
  \(J(\theta) = -\frac{1}{n}l(\theta)\)).
\item
  The log-loss function heavily penalizes confident wrong predictions.
\end{itemize}

\subsubsection{Optimizing the
parameters}\label{optimizing-the-parameters}

By minimizing the cost function \(J(\theta)\), we aim to find the
parameters \(\theta\) that maximize the likelihood of observing the
training data given the model parameters.

To achieve this, we use an iterative optimization method,
\textbf{gradient descent}, to find the values of \(\theta\) that
minimize the cost function over the training set:

\[\underset{\theta}{minimize}(J(\theta))\]

\subsubsection{Parameter estimation
methods}\label{parameter-estimation-methods}

\paragraph{Gradient Descent}\label{gradient-descent}

The gradient descent only uses the \textbf{gradient} (first derivative)
of the cost function to update the parameters \(\theta\) in the opposite
direction of the gradient, scaled by a learning rate \(\alpha\).

To compute the gradient \(\nabla J(\theta)\) we start by transforming
the expression of:

\[\log(h_\theta(x_i)) = \log\left(\frac{1}{1 + \exp(-\theta^Tx_i)}\right) = -\log(1 + \exp(-\theta^Tx_i))\]

And:

\[\log(1 - h_\theta(x_i)) = \log\left(1 - \frac{1}{1 + \exp(-\theta^Tx_i)}\right)\]

\[\log(1 - h_\theta(x_i)) = \log\left(\frac{1 + \exp(-\theta^Tx_i) - 1}{1 + \exp(-\theta^Tx_i)}\right)\]

\[\log(1 - h_\theta(x_i)) = \log\left(\frac{\exp(-\theta^Tx_i)}{1 + \exp(-\theta^Tx_i)}\right)\]

\[\log(1 - h_\theta(x_i)) = \log(\exp(-\theta^Tx_i)) - \log({1 + \exp(-\theta^Tx_i)})\]

\[\log(1 - h_\theta(x_i)) = -\theta^Tx_i - \log({1 + \exp(-\theta^Tx_i)})\]

We integrate these modifications:

\[J(\theta) = -\frac{1}{n} \sum[y_i (-\log(1 + \exp(-\theta^Tx_i))) + (1 - y_i) (-\theta^Tx_i - \log({1 + \exp(-\theta^Tx_i)}))]\]

\[J(\theta) = -\frac{1}{n} \sum[y_i (-\log(1 + \exp(-\theta^Tx_i))) + (1 - y_i) (-\theta^Tx_i - \log(1 + exp(-\theta^Tx_i)))]\]

\[J(\theta) = -\frac{1}{n} \sum[y_i (-\log(1 + \exp(-\theta^Tx_i))) -\theta^Tx_i - \log(1 + \exp(-\theta^Tx_i)) + y_i \theta^Tx_i  + y_i \log(1 + \exp(-\theta^Tx_i))]\]

\[J(\theta) = -\frac{1}{n} \sum[- y_i \log(1 + \exp(-\theta^Tx_i)) -\theta^Tx_i - \log({1 + \exp(-\theta^Tx_i)}) + y_i \theta^Tx_i  + y_i \log(1 + \exp(-\theta^Tx_i))]\]

\[J(\theta) = -\frac{1}{n} \sum[{- y_i \log(1 + \exp(-\theta^Tx_i))} -\theta^Tx_i - \log({1 + \exp(-\theta^Tx_i)}) + y_i \theta^Tx_i  + y_i \log(1 + \exp(-\theta^Tx_i))]\]

\[J(\theta) = -\frac{1}{n} \sum[-\theta^Tx_i - \log({1 + \exp(-\theta^Tx_i)}) + y_i \theta^Tx_i  ]\]

\[J(\theta) = -\frac{1}{n} \sum[y_i \theta^Tx_i  -\theta^Tx_i - \log({1 + \exp(-\theta^Tx_i)}) ]\]

with:

\[-\theta^Tx_i - \log({1 + \exp(-\theta^Tx_i)}) = - \log(\exp(\theta^T x_i)) - \log(1 + \exp(-\theta^Tx_i))\]

\[-\theta^Tx_i - \log({1 + \exp(-\theta^Tx_i)}) = -(\log(\exp(\theta^T x_i)) + \log(1 + \exp(-\theta^Tx_i)))\]

\[-\theta^Tx_i - \log({1 + \exp(-\theta^Tx_i)}) = -\log[\exp(\theta^T x_i)(1 + \exp(-\theta^Tx_i))]\]

\[-\theta^Tx_i - \log({1 + \exp(-\theta^Tx_i)}) = -\log(\exp(\theta^T x_i) + 1)\]

\[J(\theta) = -\frac{1}{n} \sum[y_i \theta^Tx_i  -\log(\exp(\theta^T x_i + 1)) ]\]

\[J(\theta) = -\frac{1}{n} \sum[y_i \theta^Tx_i  -\log(1 + \exp(\theta^T x_i)) ]\]

\[\frac{\partial}{\partial \theta_j}J(\theta) = -\frac{1}{n} \sum[y_i \frac{\partial}{\partial \theta_j} (\theta^Tx_i)  - \frac{\partial}{\partial \theta_j}\log(1 + \exp(\theta^T x_i)) ]\]

Knowing that:

\[\theta^Tx_i = \theta_1 {x_i}^{(1)} + \theta_2 {x_i}^{(2)} + \ldots + \theta_k {x_i}^{(k)}\]

\[\frac{\partial}{\partial \theta_j} (\theta^Tx_i) = x_i^{(j)}\]
\[\dfrac{\partial}{\partial \theta_j}\left(\log(1 + \exp(\theta^T x_i))\right) \underset{\log(u)^{'} = \dfrac{u^{'}}{u}} = \dfrac{\dfrac{\partial}{\partial \theta_j}(1 + \exp(\theta^T x_i))} {1 + \exp(\theta x_i)}\]

And:

\[\dfrac{\partial}{\partial \theta_j}\left(\log(1 + \exp(\theta^T x_i))\right) \underset{\log(u)^{'} = \dfrac{u^{'}}{u}} = \dfrac{\dfrac{\partial}{\partial \theta_j}(1 + \exp(\theta^T x_i))} {1 + \exp(\theta x_i)}\]
\[\dfrac{\dfrac{\partial}{\partial \theta_j}(1 + \exp(\theta^T x_i))} {1 + \exp(\theta x_i)} = \dfrac{\dfrac{\partial}{\partial \theta_j}(\exp(\theta^T x_i))} {1 + \exp(\theta x_i)}\]
\[\dfrac{\dfrac{\partial}{\partial \theta_j}(\exp(\theta^T x_i))} {1 + \exp(\theta x_i)} \underset{\exp(u)^{'} = u^{'}\exp(u)} =  \dfrac{\dfrac{\partial}{\partial \theta_j}(\theta^T x_i) * (\exp(\theta^T x_i))} {1 + \exp(\theta x_i)}\]
\[\dfrac{\dfrac{\partial}{\partial \theta_j}(\theta^T x_i) * (\exp(\theta^T x_i))} {1 + \exp(\theta x_i)} = \frac{x_i^{(j)} * (\exp(\theta^T x_i))} {1 + \exp(\theta x_i)}\]
\[\dfrac{x_i^{(j)} * (\exp(\theta^T x_i))} {1 + \exp(\theta x_i)} = x_i^{(j)} * h_\theta(x_i)\]
\[\dfrac{\partial}{\partial \theta_j}J(\theta) = -\dfrac{1}{n} \sum[y_i x_i^{(j)}  - x_i^{(j)} h_\theta(x_i)]\]

\[-\dfrac{1}{n} \sum[y_i x_i^{(j)}  - x_i^{(j)} h_\theta(x_i)] = -\dfrac{1}{n} \sum[y_i - h_\theta(x_i) ] x_i^{(j)}\]

\[\frac{\partial}{\partial \theta_j}J(\theta) = \frac{1}{n} \sum[h_\theta(x_i) - y_i ] x_i^{(j)}\]

We know that the expression of the Gradient descent to update the
weights is for the weight \(\theta_j\) :

\[\theta_j = \theta_j - \alpha \frac{\partial}{\partial \theta_j}J(\theta)\]

\[\theta_j = \theta_j - \frac{\alpha}{n} \sum[h_\theta(x_i) - y_i ] x_i^{(j)}\]

\paragraph{The Newton-Raphson algorithm}

Others algorithms can be used to find the coefficients of the logistic
regression. The Newton-Raphson algorithm is an alternative to the
gradient descent.

The Newton-Raphson method uses both the first derivative (gradient) and
the second derivative (Hessian) of the cost function. It approximates
the cost function as a quadratic form and finds the minimum of this
approximation.

\textbf{Update rule}

\[\theta = \theta - H^{-1} \nabla J(\theta)\]

Logistic regression is a regression model used to predict the
probability of a binary event based on one or more predictive variables.
In logistic regression, the likelihood function is convex and can be
maximized using the Newton-Raphson algorithm.

The Newton-Raphson algorithm is an iterative method for finding the
maximum of a function using its first and second derivatives. For
logistic regression, the likelihood function is given by:

\(L(\theta | X, y) = \prod(P(yi | x_i, \theta)^{yi} (1 - P(y_i | x_i, \theta))^{(1 - y_i)})\)

where:

\begin{itemize}
\tightlist
\item
  \(\theta\) is the vector of logistic regression coefficients,
\item
  \(X\) is the matrix of predictive variables,
\item
  \(y\) is the vector of binary response variables, and
\item
  \(P(y_i | x_i, \theta)\) is the predicted probability of the binary
  event for observation \(i\).
\end{itemize}

To maximize the likelihood function, the Newton-Raphson algorithm
updates the coefficient vector \(\theta\) at each iteration using the
following formula:

\(\theta_{i+1} = \theta_i - H^{-1} . g\)

where

\begin{itemize}
\item
  \(H = \dfrac{\partial^2L}{\partial \theta \partial\theta'}\) is the
  Hessian matrix of the likelihood function.
\item
  \(g = \dfrac{\partial L}{\partial \theta }\) is the gradient vector of
  the likelihood function, and,
\item
  \(\theta_i\) is the coefficient vector at iteration \(i\).
\end{itemize}

The Hessian matrix and gradient vector of the likelihood function are
computed using the partial derivatives of the likelihood with respect to
the coefficients \(\theta\).

Thus, in logistic regression, the Newton-Raphson algorithm is used to
estimate the coefficients by maximizing the likelihood function. This
allows the prediction of binary event probabilities based on the
predictive variables.

\textbf{Note}: The iterations stop when the difference between two
successive solution vectors becomes negligible.

\subsubsection{Convexity}\label{convexity}

Convexity is a crucial property in optimization, as it ensures that any
local minimum is also a global minimum. This makes it easier to find the
optimal solution using methods like gradient descent.

The Log Loss function is convex due to its logarithmic form, which is
always convex for positive values. While convexity guarantees that a
stationary point (where the derivative is zero) is a global minimum, it
does not guarantee the existence of such a point.

To prove the existence of a minimum, the function must also be bounded
below and attain this lower bound. For Log Loss, the function is bounded
below by zero (since the logarithm of a positive number is always
defined) and reaches this bound when predictions are perfectly accurate
(i.e., the predicted probability for the correct class is 1).

Combining convexity with the fact that Log Loss is bounded below and
attains its lower bound, we conclude that the function reaches a global
minimum when predictions are perfectly correct.

\subsection{Coefficients
significativity}\label{coefficients-significativity}

The Wald statistic allows to test the coefficients significativity
\(\hat{w_j}\). Wald statistic is given by::

\((\frac{\hat{w_j}}{\sigma(\hat{w_j})})^2\)

The added-value of the variable \(X_j\) is only real if the Wald
statistic \textgreater{} 4 \((3.84 = 1.96^2)\)

\(Wald > 4\)

\(\iff (\frac{\hat{w_j}}{\sigma(\hat{w_j})})^2 > 4\)

\(\iff \frac{\hat{w_j}}{\sigma(\hat{w_j})} > 2\)



\(\iff \hat{w_j} - 2\sigma(\hat{w_j}) > 0\)

\(\iff \hat{w_j}\) se trouve à plus de 2 écarts-type de 0


\(\iff \hat{w_j}\) se trouve à plus de 2 écarts-type de 0

\textbackslash iff  l\textquotesingle intervalle de confiance de
\(\hat{w_j}\) ne contient pas 0 à 95%

CQFD

\subsection{Model quality mesure
(Deviance)}\label{model-quality-mesure-deviance}

Cf. S.Tufféry p.315

\(n:\) number of observations\\
\(k:\) number of features

\(L(\omega_k)\) Likelihood of the "modèle ajusté"

\(L(\omega_0)\) Likelihood of the "modèle réduit à la constante"

\(L(\omega_{max})\) Likelihood of the "modèle saturé". The one the model
will compare.

The Deviance formula:

\(D(\omega_k) = -2[log(L(\omega_k)) - log(L(\omega_{max}))]\) \(^{(*)}\)

As the target is 0 or 1
\(\Longrightarrow L(\omega_{max})=1 \Longrightarrow log(L(\omega_{max}))=0\)

\(\Longrightarrow D(\omega_k) = -2[log(L(\omega_k))]\)

(*) \(D(\omega_k) = (\frac{log(L(\omega_k))}{log(L(\omega_{max}))}^2)\)

The goal of the logistic regression is to maximise the Likelihood which
is equivalent to minimize the Deviance.

The Deviance is equivalent to the SCE for the linear regression.

\end{document}